\chapter{2011年湖南卷（理）}

\section{单选题}

\begin{problem}
若 \(a,b \in \R\)，\(\i\) 为虚数单位，且 \((a + \i)\i = b + \i\)，则（\quad）

\choices
    {\(a = 1\)，\(b = 1\)}
    {\(a = -1\)，\(b = 1\)}
    {\(a = -1\)，\(b = -1\)}
    {\(a = 1\)，\(b = -1\)}
\end{problem}

\begin{answer}
D
\end{answer}

\begin{problem}
设 \(M = \{1, 2\}\)，\(N = \{a^2\}\)，则 “\(a = 1\)” 是 “\(N \subseteq M\)” 的（\quad）

\choices
    {充分不必要条件}
    {必要不充分条件}
    {充分必要条件}
    {既不充分又不必要条件}
\end{problem}

\begin{answer}
A
\end{answer}

\begin{problem}
如图是某几何体的三视图，则该几何体的体积为（\quad）

\bitmapfigure[max width=0.6\linewidth]{img/2011/hunan_science/q3_fig1.png}

\choices
    {\(\frac{9}{2}\pi + 12\)}
    {\(\frac{9}{2}\pi + 18\)}
    {\(9\pi + 42\)}
    {\(36\pi + 18\)}
\end{problem}

\begin{answer}
B
\end{answer}

\begin{problem}
通过随机询问 \(110\text{ 名}\) 性别不同的大学生是否爱好某项运动，得到如下的列联表：

\begin{center}
\begin{tabular}{c|ccc}
 & 男 & 女 & 总计 \\
\hline
爱好 & 40 & 20 & 60 \\
不爱好 & 20 & 30 & 50 \\
总计 & 60 & 50 & 110
\end{tabular}
\end{center}

由 \(K^2 = \frac{n(ad - bc)^2}{(a + b)(c + d)(a + c)(b + d)}\) 算得 \(K^{2} = \frac{110 \times (40 \times 30 - 20 \times 20)^{2}}{60 \times 50 \times 60 \times 50} \approx 7.8\). 附表：

\begin{center}
\begin{tabular}{c|ccc}
\(P(K^2\geqslant k)\) & 0.050 & 0.010 & 0.001 \\
\hline
\(k\) & 3.841 & 6.635 & 10.828
\end{tabular}
\end{center}

参照附表，得到的正确结论是（\quad）

\choices
    {在犯错误的概率不超过 \(0.1\%\) 的前提下，认为“爱好该项运动与性别有关”}
    {在犯错误的概率不超过 \(0.1\%\) 的前提下，认为“爱好该项运动与性别无关”}
    {有 \(99\%\) 以上的把握认为“爱好该项运动与性别有关”}
    {有 \(99\%\) 以上的把握认为“爱好该项运动与性别无关”}
\end{problem}

\begin{answer}
C
\end{answer}

\begin{problem}
设双曲线 \(\frac{x^2}{a^2} - \frac{y^2}{9} = 1\)（\(a > 0\)）的渐近线方程为 \(3x \pm 2y = 0\)，则 \(a\) 的值为（\quad）

\choices
    {4}
    {3}
    {2}
    {1}
\end{problem}

\begin{answer}
C
\end{answer}

\begin{problem}
由直线 \(x = -\frac{\pi}{3}\)，\(x = \frac{\pi}{3}\)，\(y = 0\) 与曲线 \(y = \cos x\) 所围成的封闭图形的面积为（\quad）

\choices
    {\(\frac{1}{2}\)}
    {1}
    {\(\frac{\sqrt{3}}{2}\)}
    {\(\sqrt{3}\)}
\end{problem}

\begin{answer}
D
\end{answer}

\begin{problem}
设 \(m > 1\)，在约束条件 \(\begin{cases}
y \geqslant x,\\
y \leqslant mx,\\
x + y \leqslant 1
\end{cases}\) 下，目标函数 \(z = x + my\) 的最大值小于 \(2\)，则 \(m\) 的取值范围为（\quad）

\choices
    {\((1, 1 + \sqrt{2})\)}
    {\((1 + \sqrt{2}, +\infty)\)}
    {\((1,3)\)}
    {\((3, +\infty)\)}
\end{problem}

\begin{answer}
A
\end{answer}

\begin{problem}
设直线 \(x = t\) 与函数 \(f(x) = x^2\)，\(g(x) = \ln x\) 的图象分别交于点 \(M\)，\(N\)，则当 \(|MN|\) 达到最小时 \(t\) 的值为（\quad）

\choices
    {1}
    {\( \frac{1}{2} \)}
    {\(\frac{\sqrt{5}}{2}\)}
    {\(\frac{\sqrt{2}}{2}\)}
\end{problem}

\begin{answer}
D
\end{answer}

\section{填空题}

\begin{problem}
在直角坐标系 \(xOy\) 中，曲线 \(C_1\) 的参数方程为 \(\begin{cases}
x = \cos \alpha,\\
y = 1 + \sin \alpha
\end{cases}\) （\(\alpha\) 为参数），在极坐标系（与直角坐标系 \(xOy\) 取相同的长度单位，且以原点 \(O\) 为极点，以 \(x\) 轴正半轴为极轴）中，曲线 \(C_2\) 的方程为 \(\rho (\cos \theta - \sin \theta) + 1 = 0\)，则 \(C_1\) 与 \(C_2\) 的交点个数为\fillinblank{}.
\end{problem}

\begin{answer}
2
\end{answer}

\begin{problem}
设 \(x,y \in \R\)，且 \(xy \neq 0\)，则 \(\left(x^2 + \frac{1}{y^2}\right)\left(\frac{1}{x^2} + 4y^2\right)\) 的最小值为\fillinblank{}.
\end{problem}

\begin{answer}
9
\end{answer}

\begin{problem}
如图，\(A\)，\(E\) 是半圆周上的两个三等分点，直径 \(BC=4\)，\(AD\perp BC\)，垂足为 \(D\)，\(BE\) 与 \(AD\) 相交于点 \(F\)，则 \(AF\) 的长为\fillinblank{}.

\bitmapfigure[max width=0.34\linewidth]{img/2011/hunan_science/q11_fig1.png}
\end{problem}

\begin{answer}
\(\frac{2\sqrt{3}}{3}\)
\end{answer}

\begin{problem}
设 \(S_{n}\) 是等差数列 \(\{a_{n}\}\)（\(n \in \mathbf{N}^{*}\)）的前 \(n\) 项和，且 \(a_{1} = 1\)，\(a_{4} = 7\)，则 \(S_{5} =\)\fillinblank{}.
\end{problem}

\begin{answer}
25
\end{answer}

\begin{problem}
若执行如图所示的框图，输入 \(x_{1} = 1\)，\(x_{2} = 2\)，\(x_{3} = 3\)，\(\overline{x} = 2\)，则输出的数等于\fillinblank{}.

\bitmapfigure[max width=0.32\linewidth]{img/2011/hunan_science/q13_fig1.png}
\end{problem}

\begin{answer}
\(\frac{2}{3}\)
\end{answer}

\begin{problem}
在边长为 1 的正三角形 \(ABC\) 中，设 \(\overrightarrow{BC} = 2\overrightarrow{BD}\)，\(\overrightarrow{CA} = 3\overrightarrow{CE}\)，则 \(\overrightarrow{AD} \cdot \overrightarrow{BE} =\)\fillinblank{}.
\end{problem}

\begin{answer}
\(-\frac{1}{4}\)
\end{answer}

\begin{problem}
如图所示，\(EFGH\) 是以 \(O\) 为圆心，半径为 1 的圆的内接正方形，将一颗豆子随机地扔到该圆内，用 \(A\) 表示事件“豆子落在正方形 \(EFGH\) 内”，\(B\) 表示事件“豆子落在扇形 \(OHE\)（阴影部分）内”，则
\begin{enumerate}
\item \(P(A)=\)\fillinblank{}；
\item \(P(B|A)=\)\fillinblank{}.
\end{enumerate}

\bitmapfigure[max width=0.35\linewidth]{img/2011/hunan_science/q15_fig1.png}
\end{problem}

\begin{answer}
\(\frac{2}{\pi}\)，\(\frac{1}{4}\)
\end{answer}

\begin{problem}
对于 \(n \in \mathbf{N}^*\)，将 \(n\) 表示为 \(n = a_0 \times 2^k + a_1 \times 2^{k-1} + a_2 \times 2^{k-2} + \cdots + a_{k-1} \times 2^1 + a_k \times 2^0\). 当 \(i = 0\) 时，\(a_i = 1\)，当 \(1 \leqslant i \leqslant k\) 时，\(a_i\) 为 0 或 1. 记 \(I(n)\) 为上述表示中 \(a_i\) 为 0 的个数 （例如： \(1 = 1 \times 2^0\)，\(4 = 1 \times 2^2 + 0 \times 2^1 + 0 \times 2^0\)，则 \(I(1) = 0\)，\(I(4) = 2\)），则
\begin{enumerate}
    \item \(I(12) =\)\fillinblank{}；
    \item \(\displaystyle\sum_{n=1}^{127} 2^{I(n)}=\)\fillinblank{}.
\end{enumerate}
\end{problem}

\begin{answer}
2，1093
\end{answer}

\section{解答题}

\begin{problem}
在 \(\triangle ABC\) 中，角 \(A\)，\(B\)，\(C\) 所对的边分别为 \(a\)，\(b\)，\(c\). 且满足 \(c \sin A = a \cos C\).
\begin{enumerate}
\item 求角 \(C\) 的大小；
\item 求 \(\sqrt{3} \sin A - \cos \left(B + \frac{\pi}{4}\right)\) 的最大值，并求取得最大值时角 A，B 的大小.
\end{enumerate}
\end{problem}

\begin{solution}
由正弦定理，存在正数 \(R\)，使得 \(a=2R\sin A\)，\(c=2R\sin C\). 因为 \(A\) 是三角形内角，\(\sin A>0\)，所以条件等价于
\[
2R\sin C\sin A=2R\sin A\cos C\quad\Longleftrightarrow\quad \sin C=\cos C.
\]
又 \(0<C<\pi\)，故 \(C=\frac{\pi}{4}\). 从而 \(A+B=\frac{3\pi}{4}\)，且 \(0<A<\frac{3\pi}{4}\).

由 \(B=\frac{3\pi}{4}-A\)，有 \(B+\frac{\pi}{4}=\pi-A\)，因此
\[
\sqrt{3}\sin A-\cos\left(B+\frac{\pi}{4}\right)=\sqrt{3}\sin A+\cos A=2\cos\left(A-\frac{\pi}{3}\right)\leqslant 2.
\]
当且仅当 \(A=\frac{\pi}{3}\) 时等号成立，此时 \(B=\frac{5\pi}{12}\). 所以最大值为 \(2\)，取得最大值时 \(A=\frac{\pi}{3}\)，\(B=\frac{5\pi}{12}\).
\end{solution}

\begin{problem}
某商店试销某种商品 \(20\text{ 天}\)，获得如下数据：

\begin{center}
\begin{tabular}{c|cccc}
日销售量（件） & 0 & 1 & 2 & 3 \\
\hline
频数 & 1 & 5 & 9 & 5
\end{tabular}
\end{center}

试销结束后（假设该商品的日销售量的分布规律不变），设某天开始营业时有该商品 \(3\text{ 件}\)，当天营业结束后检查存货，若发现存货少于 \(2\text{ 件}\)，则当天进货补充至 \(3\text{ 件}\)，否则不进货，将频率视为概率.

\begin{enumerate}
\item 求当天商品不进货的概率；
\item 记 \(X\) 为第二天开始营业时该商品的件数，求 \(X\) 的分布列和数学期望.
\end{enumerate}
\end{problem}

\begin{solution}
设当天的日销售量为 \(Y\). 当天开始营业时有 \(3\) 件商品，若当天销售量为 \(0\) 件或 \(1\) 件，营业结束时存货分别为 \(3\) 件或 \(2\) 件，不需要进货；若销售量为 \(2\) 件或 \(3\) 件，营业结束时存货少于 \(2\) 件，需要进货.

\begin{enumerate}
\item 当天商品不进货的概率为
\[
P(Y=0)+P(Y=1)=\frac{1+5}{20}=\frac{3}{10}.
\]

\item 第二天开始营业时，若 \(Y=1\)，存货为 \(2\) 件；若 \(Y=0\)，存货为 \(3\) 件；若 \(Y=2\) 或 \(Y=3\)，当天进货后存货补充至 \(3\) 件. 因此 \(X\) 只取 \(2\)，\(3\)，且
\[
P(X=2)=P(Y=1)=\frac{5}{20}=\frac{1}{4},
\]
\[
P(X=3)=P(Y=0)+P(Y=2)+P(Y=3)=\frac{1+9+5}{20}=\frac{3}{4}.
\]
分布列为
\begin{center}
\begin{tabular}{c|cc}
\(X\) & 2 & 3 \\
\hline
\(P\) & \(\frac{1}{4}\) & \(\frac{3}{4}\)
\end{tabular}
\end{center}
所以
\[
E(X)=2\times\frac{1}{4}+3\times\frac{3}{4}=\frac{11}{4}.
\]
\end{enumerate}
\end{solution}

\begin{problem}
如图，在圆锥 \(PO\) 中，已知 \(PO = \sqrt{2}\)，\(\odot O\) 的直径 \(AB = 2\)，\(C\) 是 \(\widehat{AB}\) 的中点，\(D\) 为 \(AC\) 的中点.
\begin{enumerate}
\item 证明：平面 \(POD \perp\) 平面 \(PAC\)；
\item 求二面角 \(B - PA - C\) 的余弦值.
\end{enumerate}

\bitmapfigure[max width=0.36\linewidth]{img/2011/hunan_science/q19_fig1.png}
\end{problem}

\begin{solution}
不妨在底面所在平面内建立直角坐标系，取
\(O=(0,0,0)\)，\(P=(0,0,\sqrt{2})\)，\(A=(-1,0,0)\)，\(B=(1,0,0)\)，\(C=(0,1,0)\).
于是 \(D=\left(-\frac{1}{2},\frac{1}{2},0\right)\). 取 \(C\) 在另一侧时只是坐标方向改变，结论不变.

\begin{enumerate}
\item 因为 \(D\) 是圆心为 \(O\) 的圆内弦 \(AC\) 的中点，所以 \(OD\perp AC\). 又因为 \(PO\) 垂直于底面，故 \(PO\perp AC\). 直线 \(PO\)，\(OD\) 是平面 \(POD\) 内相交的两条直线，且都垂直于 \(AC\)，所以 \(AC\perp\) 平面 \(POD\). 由于 \(AC\) 在平面 \(PAC\) 内，故平面 \(POD\perp\) 平面 \(PAC\).

\item 平面 \(PAB\) 的方程为 \(y=0\)，可取其一个法向量为 \(\bs{n}_1=(0,1,0)\). 在平面 \(PAC\) 内取
\(\overrightarrow{AP}=(1,0,\sqrt{2})\)，\(\overrightarrow{AC}=(1,1,0)\)，
从而可取平面 \(PAC\) 的法向量为
\[
\bs{n}_2=\overrightarrow{AP}\times\overrightarrow{AC}=(-\sqrt{2},\sqrt{2},1).
\]
设二面角 \(B-PA-C\) 的平面角为 \(\theta\)，则它是平面 \(PAB\) 与平面 \(PAC\) 的锐角，故
\[
\cos\theta=\frac{|\bs{n}_1\cdot\bs{n}_2|}{|\bs{n}_1||\bs{n}_2|}=\frac{\sqrt{2}}{\sqrt{5}}=\frac{\sqrt{10}}{5}.
\]
所以二面角 \(B-PA-C\) 的余弦值为 \(\frac{\sqrt{10}}{5}\).
\end{enumerate}
\end{solution}

\begin{problem}
如图，长方体物体 \(E\) 在雨中沿面 \(P\)（面积为 \(S\)）的垂直方向作匀速移动，速度为 \(v\)（\(v > 0\)），雨速沿 \(E\) 移动方向的分速度为 \(c\)（\(c \in \R\)）. \(E\) 移动时单位时间内的淋雨量包括两部分：
\begin{circlelist}
\item \(P\) 或 \(P\) 的平行面 （只有一个面淋雨） 的淋雨量，假设其值与 \(|v - c| \times S\) 成正比，比例系数为 \(\frac{1}{10}\)；
\item 其它面的淋雨量之和，其值为 \(\frac{1}{2}\)；
\end{circlelist}
记 \(y\) 为 \(E\) 移动过程中的总淋雨量，当移动距离 \(d = 100\)，面积 \(S = \frac{3}{2}\) 时.
\begin{enumerate}
\item 写出 \(y\) 的表达式；
\item 设 \(0 < v \leqslant 10\)，\(0 < c \leqslant 5\)，试根据 \(c\) 的不同取值范围，确定移动速度 \(v\)，使 \(y\) 最少.
\end{enumerate}

\problemresetlayout
\FigureLayoutDeclare{img/2011/hunan_science/q20_fig1.png}{5.3cm}{33bp 25bp 37bp 30bp}
\bitmapfigure[width=0.35\linewidth]{img/2011/hunan_science/q20_fig1.png}
\end{problem}

\begin{solution}
\begin{enumerate}
\item 单位时间内，面 \(P\) 或其平行面的淋雨量为 \(\frac{1}{10}|v-c|S\)，其它面的淋雨量之和为 \(\frac{1}{2}\)，所以单位时间内的总淋雨量为
\[
\frac{1}{10}|v-c|S+\frac{1}{2}.
\]
物体移动距离为 \(100\)，速度为 \(v\)，移动时间为 \(\frac{100}{v}\). 代入 \(S=\frac{3}{2}\)，得
\[
y=\frac{100}{v}\left(\frac{1}{10}|v-c|\cdot\frac{3}{2}+\frac{1}{2}\right)=\frac{15|v-c|+50}{v}.
\]

\item 当 \(0<v<c\) 时，
\[
y=\frac{50+15c}{v}-15,
\]
该式随 \(v\) 增大而减小；当 \(c\leqslant v\leqslant 10\) 时，
\[
y=15+\frac{50-15c}{v}.
\]

当 \(0<c<\frac{10}{3}\) 时，\(50-15c>0\)，后一个式子随 \(v\) 增大而减小，故总淋雨量在 \(v=10\) 时最少.

当 \(c=\frac{10}{3}\) 时，后一个式子恒为 \(15\)，故所有 \(v\in\left[\frac{10}{3},10\right]\) 都使总淋雨量最少.

当 \(\frac{10}{3}<c\leqslant 5\) 时，\(50-15c<0\)，后一个式子随 \(v\) 增大而增大，且前一个式子在 \(0<v<c\) 时随 \(v\) 增大而减小，所以总淋雨量在 \(v=c\) 时最少.
\end{enumerate}
\end{solution}

\begin{problem}
如图，椭圆 \(C_{1}:\frac{x^{2}}{a^{2}} +\frac{y^{2}}{b^{2}} = 1\)（\(a > b > 0\)）的离心率为 \(\frac{\sqrt{3}}{2}\)，\(x\) 轴被曲线 \(C_2\)：\(y = x^{2} - b\) 截得的线段长等于 \(C_1\) 的长半轴长.
\begin{enumerate}
\item 求 \(C_1\)，\(C_2\) 的方程；
\item 设 \(C_2\) 与 \(y\) 轴的交点为 \(M\)，过坐标原点 \(O\) 的直线 \(l\) 与 \(C_2\) 相交于点 \(A\)，B，直线 \(MA\)，\(MB\) 分别与 \(C_1\) 相交于点 \(D\)，\(E\).
\begin{enumerate}
\item 证明：\(MD \perp ME\)；
\item 记 \(\triangle MAB\)，\(\triangle MDE\) 的面积分别是 \(S_{1}\)，\(S_{2}\). 问：是否存在直线 \(l\)，使得 \(\frac{S_{1}}{S_{2}} = \frac{17}{32}\)，请说明理由.
\end{enumerate}
\end{enumerate}

\bitmapfigure[max width=0.42\linewidth]{img/2011/hunan_science/q21_fig1.png}
\end{problem}

\begin{solution}
\begin{enumerate}
\item 设椭圆的半焦距为 \(c_0\). 由离心率条件 \(\frac{c_0}{a}=\frac{\sqrt{3}}{2}\) 和 \(c_0^2=a^2-b^2\)，
得 \(b=\frac{a}{2}\). 曲线 \(C_2\) 与 \(x\) 轴的交点满足 \(x^2-b=0\)，所以截得线段的两个端点横坐标为 \(\pm\sqrt b\)，线段长为 \(2\sqrt b\). 由题意 \(2\sqrt b=a\)，即 \(a^2=4b\). 联立 \(b=\frac{a}{2}\) 和 \(a^2=4b\)，并由 \(a>0\) 得 \(a=2\)，\(b=1\). 因此 \(C_1:\frac{x^2}{4}+y^2=1\)，\(C_2:y=x^2-1\).

\item 设直线 \(l\) 的斜率为 \(k\)，则 \(l:y=kx\). 竖直直线与 \(C_2\) 至多有一个交点，故能够与 \(C_2\) 交于两点的直线 \(l\) 确有斜率. 设 \(A\)，\(B\) 的横坐标分别为 \(t\)，\(s\)，则
\(t^2-kt-1=0\)，\(t+s=k\)，\(ts=-1\).
又 \(M=(0,-1)\)，故
\(A=(t,t^2-1)\)，\(B=(s,s^2-1)\).

\begin{enumerate}
\item 设 \(T=(u,u^2-1)\) 是 \(C_2\) 上一点，直线 \(MT\) 上的点可表示为
\[
M+\lambda(T-M)=\left(\lambda u,-1+\lambda u^2\right).
\]
将其代入 \(C_1\) 的方程，得
\[
\frac{\lambda^2u^2}{4}+\left(-1+\lambda u^2\right)^2=1.
\]
除去对应于 \(M\) 的根 \(\lambda=0\)，另一个根满足
\[
\lambda\left(\lambda\left(\frac{u^2}{4}+u^4\right)-2u^2\right)=0,
\]
所以 \(\lambda=\frac{8}{4u^2+1}\). 因此
\(\overrightarrow{MD}=\left(\frac{8t}{4t^2+1},\frac{8t^2}{4t^2+1}\right)\)，
\(\overrightarrow{ME}=\left(\frac{8s}{4s^2+1},\frac{8s^2}{4s^2+1}\right)\).
于是
\[
\overrightarrow{MD}\cdot\overrightarrow{ME}
=\frac{64ts(1+ts)}{(4t^2+1)(4s^2+1)}=0.
\]
故 \(MD\perp ME\).

\item 由行列式表示三角形面积，
\[
2S_1=\left|\det\begin{pmatrix}t&t^2\\s&s^2\end{pmatrix}\right|=|ts(s-t)|=|s-t|.
\]
同理，由上面得到的 \(\overrightarrow{MD}\)，\(\overrightarrow{ME}\)，有
\[
2S_2=\frac{64|ts||s-t|}{(4t^2+1)(4s^2+1)}=\frac{64|s-t|}{(4t^2+1)(4s^2+1)}.
\]
所以
\[
\frac{S_1}{S_2}=\frac{(4t^2+1)(4s^2+1)}{64}
=\frac{16t^2s^2+4(t^2+s^2)+1}{64}
=\frac{25+4k^2}{64},
\]
其中使用了 \(ts=-1\) 及 \(t+s=k\). 若要求该比值为 \(\frac{17}{32}\)，则
\[
\frac{25+4k^2}{64}=\frac{17}{32}\quad\Longleftrightarrow\quad k^2=\frac{9}{4}.
\]
取 \(k=\frac{3}{2}\) 或 \(k=-\frac{3}{2}\)，相应直线均能与 \(C_2\) 交于两点，故存在这样的直线 \(l\).
\end{enumerate}
\end{enumerate}
\end{solution}

\begin{problem}
已知函数 \(f(x) = x^3\)，\(g(x) = x + \sqrt{x}\).
\begin{enumerate}
\item 求函数 \( h(x) = f(x) - g(x) \) 的零点个数，并说明理由；
\item 设数列 \(\{a_{n}\}\)（\(n \in \mathbf{N}^{*}\)）满足 \(a_{1} = a\)（\(a > 0\)），\(f(a_{n+1}) = g(a_{n})\)，证明：存在常数 \(M\)，使得对于任意的 \(n \in \mathbf{N}^{*}\)，都有 \(a_{n} \leqslant M\).
\end{enumerate}
\end{problem}

\begin{solution}
\begin{enumerate}
\item 函数 \(g\) 的定义域为 \(x\geqslant0\). 令 \(u=\sqrt{x}\)，则 \(u\geqslant0\)，且
\[
h(x)=x^3-x-\sqrt{x}=u^6-u^2-u=u\left(u^5-u-1\right).
\]
记 \(\varphi(u)=u^5-u-1\). 当 \(u\geqslant0\) 时，
\[
\varphi'(u)=5u^4-1.
\]
所以 \(\varphi\) 在 \(\left(0,5^{-1/4}\right)\) 上递减，在 \(\left(5^{-1/4},+\infty\right)\) 上递增. 又 \(\varphi(0)=-1\)，且 \(\varphi(u)\) 当 \(u\) 趋于正无穷时趋于正无穷，因此 \(\varphi(u)=0\) 有且只有一个正根，记为 \(u_0\). 所以 \(h(x)=0\) 的根为 \(x=0\) 和 \(x=u_0^2\)，共有两个.

\item 由 \(a_1=a>0\) 及递推条件，所有 \(a_n>0\)，且
\[
a_{n+1}=\left(a_n+\sqrt{a_n}\right)^{1/3}.
\]
令 \(\alpha=u_0^2\)，则 \(\alpha^3=\alpha+\sqrt\alpha\). 记 \(T(x)=\left(x+\sqrt{x}\right)^{1/3}\)（\(x>0\)）.
函数 \(T\) 在正半轴上单调递增，且 \(T(\alpha)=\alpha\). 由第一问中 \(h(x)\) 的符号可知，当 \(x\geqslant\alpha\) 时 \(x^3\geqslant x+\sqrt{x}\)，从而 \(T(x)\leqslant x\). 因此取
\[
M=\max\{a,\alpha\},
\]
则 \(a_1\leqslant M\)，并且由 \(T\) 的单调性及 \(M\geqslant\alpha\)，有 \(a_{n+1}=T(a_n)\leqslant T(M)\leqslant M\). 归纳可得，对任意 \(n\in\mathbf N^*\)，都有 \(a_n\leqslant M\).
\end{enumerate}
\end{solution}
